\documentclass[11pt]{article}
\usepackage[utf8]{inputenc}
\usepackage[T1]{fontenc}
\usepackage{lmodern}
\usepackage[margin=2.5cm]{geometry}
\usepackage{fancyhdr}\pagestyle{fancy}
\fancyhead[L]{\small Computer Science HL}
\fancyhead[R]{\small A1.1.1--A1.1.9}
\fancyfoot[C]{\thepage}\setlength{\headheight}{14pt}
\usepackage{enumitem}
\usepackage{listings}
\usepackage{xcolor}
\usepackage{tabularx}
\usepackage{booktabs}
\usepackage{hyperref}
\lstset{basicstyle=\ttfamily\small, keywordstyle=\color{blue}, language=Python, frame=single}
% user environments: an inline one (font switch and a label) and a block one (wraps quote)
\newenvironment{solution}{\par\noindent\textbf{Solution.}\ \itshape}{\par}
\newenvironment{hint}{\begin{quote}\small\textbf{Hint:}\ }{\end{quote}}
\begin{document}
{\Large\bfseries Computer Hardware and Operation\par}
Henry Yang \hfill 22 September 2026

\section{The CPU}
The central processing unit (CPU) fetches, decodes and executes instructions. The \textbf{control unit} (CU) coordinates the \emph{fetch--decode--execute} cycle, and the \textbf{arithmetic logic unit} (ALU) performs calculations such as \texttt{ADD}, \texttt{SUB} and comparisons.\footnote{Some architectures add a dedicated floating-point unit.}

\begin{enumerate}[label=(\alph*)]
  \item The \textbf{program counter} holds the address of the next instruction.
  \item The \textbf{memory address register} holds the address being read or written.
  \item The \textbf{memory data register} holds the data just read or about to be written.
\end{enumerate}

Registers are described in Table~\ref{tab:regs}; see also \url{https://en.wikipedia.org/wiki/Processor_register} and the \href{https://example.org/notes}{course notes}.

\begin{table}[h]
\centering
\begin{tabularx}{\textwidth}{lX}
\toprule
Register & Purpose \\
\midrule
PC & Address of the next instruction to fetch. \\
MAR & Address of the memory cell to access. \\
MDR & Data read from or written to memory. \\
\bottomrule
\end{tabularx}
\caption{Processor registers.}\label{tab:regs}
\end{table}

\section{A short program}
The following Python program adds two numbers:
\begin{lstlisting}
def add(a, b):
    # return the sum
    return a + b

print(add(2, 3))
\end{lstlisting}
Calling \lstinline|add(2, 3)| prints \texttt{5}.

\begin{itemize}[noitemsep]
  \item \textbf{Primary memory} (RAM) is volatile.
  \item \textbf{Secondary memory} keeps data when the power is off.
  \begin{itemize}
    \item Hard disk drives use magnetic platters.
    \item Solid-state drives use flash memory.
  \end{itemize}
\end{itemize}

\begin{description}
  \item[Cache] A small, fast memory close to the CPU.
  \item[Bus] A set of wires carrying data, addresses and control signals.
\end{description}

\newpage
\section{Operating systems}
An operating system manages \textit{resources}: memory, the processor and peripherals. Its \textbf{scheduler} decides which process runs next; a \emph{context switch} saves and restores register state. See Section~\ref{sec:q}.

\section{Questions}\label{sec:q}
\begin{enumerate}[leftmargin=*]
  \item Outline the function of the MAR. \hfill [2]
  \item Describe one difference between RAM and ROM. \hfill [2]
\end{enumerate}

\begin{hint}
  Think about what happens to the contents of each kind of memory when the power is switched off, and which of them the processor can write to.
\end{hint}

\begin{solution}
The MAR holds the address of the memory location that is about to be read from or written to, so that the address bus can be driven with it while the data travels over the data bus.
\end{solution}
\end{document}
